\scanpage{10}%
toutes égales à 1\handwritten{)}. On trouve alors une description complète de la variation des valeurs propres mod $p$ de frobenius opérant sur $E(Y)$, pour $Y$ variable avec un point de $U_d$, en termes essentiellement de la seule représentation de $G$ dans un vectoriel de dimension finie sur $\underline{\mathrm F}_p$, correspondant aux $E(Y,\underline O_Y)_{ss}$ (donc essentiellement aux ``matrices de Hasse-Witt des $Y$ en dimension critique), - et bien entendu de la distribution des \deleted{frobenius} éléments de frobenius dans $G$ en termes d'un point variable de $U_d$.

8.5.10. L'homomorphisme \deleted{\unsure{$\varphi$}}fondamental $\varphi$ de (8.5.1) est \deleted{surjectif (resp.} injectif\deleted{)} \inserted{\handwritten{donc bijectif}} si et seulement si $w$ dans 8.5.\deleted{\unsure{3}}\handwritten{1} est \deleted{injectif (resp. si $w$ est} surjectif, \inserted{\handwritten{et il l'est si $({}_p\mathrm{Tors}^{n}(X))_{ss}=0$,}} le seulement si étant vrai sous la réserve $\mathrm{Ker}\,v=0$ de 8.5.7); il suffit en particulier pour ceci que $X$ n'ait pas de $p$-torsion en \deleted{caractéristique} dimension $n\deleted{\unsure{-1}}$ \deleted{(resp. en dimension \unsure{$n$}\deleted{\unsure{, ou}} $n+1$,} \handwritten{(}ou ce qui revient au même par dualité, en dimension $n\handwritten{+1}$\handwritten{)}. Lorsque $X$ se \deleted{Je pense que même si} relève (avec sa polarisation) en caractéristique nulle, ce\deleted{s}\handwritten{tte} condition\deleted{s} \deleted{s'exp} \inserted{suffisante\deleted{s}} doi\deleted{vent}\handwritten{t} aussi pouvoir s'exprimer (une fois qu'on disposera de fondements satisfaisants pour la cohomologie cristalline) par l\deleted{es}\handwritten{a} condition\deleted{s} analogue\deleted{s} sur \inserted{torsion\handwritten{ de}} la cohomologie $p$-\uline{adique} de la variété relevée (sur un corps algébriquement clos, ou si on veut, sur le corps des complexes ...): \deleted{on réécrit (8.5.3) en termes de $p$-torsion de cette variété tensorisée par $\underline Z/p\underline Z$ (mais en supprimant les indices $ss$ , qui n'ont plus de sens), et on demande que $w_p$ ainsi défini soit surjectif resp. injectif ....}

\deleted{Je pense que même si $n=\dim X=2$, il peut arriver que $\varphi$ ne soit pas surjectif, \deleted{\unsure{i.e}}et qu'on ne peut pas dans l'inégalité (8.5.6) supprimer le terme correctif $\dim\mathrm{Ker}\,w$ ou $t^{n-1}(X,p)$. En effet\deleted{\unsure{,}}, si $X$ est une}
% archive sheet number: 9

